\documentclass[11pt]{article}

\usepackage{amsmath,amssymb}
\usepackage{amsthm}
\usepackage[margin=1in]{geometry}
\usepackage{hyperref}
\usepackage[round,authoryear]{natbib}

\newtheorem{result}{Result}
\newtheorem{corollary}{Corollary}

\title{Pass-Through, Welfare, and Incidence under Product Returns\\
  in Perfect Competition}
\author{Draft notes}
\date{\today}

\begin{document}
\maketitle

%=============================================================================
\section{Introduction}
%=============================================================================

Under the European Union's standard value-added tax (VAT), the seller charges
tax on the invoice and remits it to the government, known as seller
remittance.  For many cross-border business-to-business supplies, the reverse
charge mechanism instead shifts that obligation to the buyer, known as buyer
remittance.  Under reverse charge, the seller invoices without VAT, and the
buyer accounts for the tax.  A similar distinction appears in the United States
between sales tax, collected and remitted by the seller, and use tax, which
the buyer self-assesses on out-of-state purchases when the seller has no
nexus.\footnote{Since the 2018 \emph{Wayfair} decision, nearly every state
  with a sales tax has adopted ``marketplace facilitator'' laws that require
  platforms such as Amazon, eBay, and Etsy to collect and remit tax on behalf
  of their third-party sellers.  As a result, the platform, not the
  independent seller, has become the statutory remittance agent for the vast
  majority of U.S.\ e-commerce transactions.  This shift has reduced the
  frequency with which consumers must self-assess use tax, while concentrating
  remittance responsibility in a small number of large platforms.}

When the transaction is final and the good is kept, these remittance rules
are economically equivalent.  The treasury receives the same revenue, and
which party remits is only an administrative detail.  A return changes that
picture, because the tax paid to the government must itself be unwound.
Reclaiming that tax typically requires filing and can involve delay or
incomplete recovery, so the remitter's effective tax refund is smaller than
the statutory amount.  We take the extreme case in which the remitter
recovers none of the tax from the government on a returned unit---whether
the seller or the buyer remitted.\footnote{Partial reclaim from the
  government would lie between full recovery of $\tau$ and this zero-reclaim
  benchmark.  We use the extreme to put both remittance regimes on the same
  footing with respect to treasury friction.}
Even so, the private refund between buyer and seller still differs by
regime.  Under seller remittance the seller refunds the full tax-inclusive
list price, while under buyer remittance the seller refunds only the
pre-tax price and the buyer does not recover the tax.  The keep-or-return
decision and the seller's effective margin therefore depend on which party
remits.

Most incidence and pass-through analysis abstracts from this difference.  It
assumes remittance neutrality and treats a return as undoing the purchase, so
pass-through, incidence, and deadweight loss are read off demand and supply
alone.  That abstraction is hard to justify in practice.  About one in five
online purchases is
returned,\footnote{National Retail Federation and Happy Returns, \emph{2025
  Retail Returns Landscape}, October 15, 2025,
  \url{https://nrf.com/research/2025-retail-returns-landscape}.
  The report estimates that 19.3\,\% of online sales will be returned in 2025.}
and a return is not a free reversal of the transaction.  Never buying and buying then returning leave the same final
allocation---the consumer does not keep the good---yet a returned purchase
incurs shipping, hassle, and handling that a non-purchase avoids.  Demand and
supply alone therefore omit a margin between purchase and final ownership that
matters for welfare.

What objects then replace the classical demand and supply curves, and how
misleading is tax analysis that omits returns?  We study these questions
under perfect competition with product returns and seller handling costs.
We characterize pass-through, incidence, and welfare losses with returns,
compare them to the no-return benchmark, and identify where omitting the
return margin changes conclusions.

Even with a return margin, surplus continues to rest on a small set of market
aggregates.  What changes is which aggregates matter and at which prices they
are evaluated.  Consumer surplus depends on only two aggregates: gross demand
and the return rate.  Heterogeneity in tastes, match values, and return costs
enters welfare solely through these aggregates.  Producer surplus retains the
usual geometry---area to the left of the supply curve---but is evaluated at
the effective margin, list price minus expected refunds and handling, rather
than at list price.

Because consumers pay the list price while firms respond to the effective
margin, how that margin moves with the list price governs pass-through from
cost shocks into list prices.  The response of the effective margin to the
list price reflects two forces.  First, a higher price directly raises
revenue on kept units.  Second, the return rate itself may change with price.
When the behavioral response dominates---consumers become more likely to return
as price rises---the effective margin can actually fall with the list price.
This creates a downward-sloping supply curve in list-price space.  Because the
refund after a return differs by remittance regime, seller and buyer remittance
generate different effective margins and equilibria, and therefore different
pass-through and incidence.

We show that pass-through of a one-dollar production-cost shock equals that
of a seller-remitted unit tax, but differs under buyer remittance.  The ratio
of handling to production-cost pass-through equals the aggregate return rate.
When the behavioral response of returns to price is strong enough that the
effective margin falls in the list price, buyer- and seller-remitted
pass-through carry opposite signs---a reversal of the classical tax-neutrality
intuition.  Which regime yields a negative pass-through depends on the
relative magnitudes of the supply and demand elasticities.  Intuitively,
negative pass-through arises when the list-price supply curve slopes
downward---a higher price lowers the effective margin and hence quantity
supplied---and more steeply than demand.  In this region, a cost increase
creates excess demand, and the equilibrium price must fall so that supply
expands by more than demand.

Despite pass-through of a cost shock potentially being negative, the
incidence ratio---the burden on consumers relative to producers---is always
positive and identical across production cost, handling cost, and
seller-remitted taxes.  For those shocks, consumer and producer surplus
therefore always move together, as in the classical model without returns.
Buyer remittance adds a direct loss to returners at fixed $p$, so its
incidence ratio differs and need not inherit that common value.

The common direction of those price-only surpluses nevertheless hinges on the
sign of pass-through and on remittance.  When cost pass-through is negative,
production cost, handling cost, and seller tax shocks raise both surpluses,
yielding a Pareto improvement, while the buyer-remitted tax need not.
In contrast, when pass-through is positive, the price-only shocks reduce both
surpluses.  Remittance therefore matters for the welfare direction of a tax
when pass-through is negative.

These welfare effects also reshape first-order changes in total private
surplus.  We decompose the change from a cost or tax increase into an
ex-ante piece on the purchase margin and an ex-post piece on the
keep-or-return margin.  For a production-cost or seller-tax shock, the
ex-ante piece is the classical rectangle from fewer units sold.  The
ex-post piece captures how the
price change reallocates units between kept and returned states, with each
additional return wasting the list price plus handling cost.  Without returns, private surplus falls by exactly the market quantity.
With returns, that fall is joined by an ex-post adjustment that the
classical no-return formula omits entirely.  The same two pieces
appear for the other shocks after a simple rescaling.  A handling-cost shock
scales both by the return rate, and under buyer remittance both are scaled by
the response of the effective margin to list price.  For taxes we then add
treasury revenue $G=\tau Q(\tau)$ back to obtain social surplus, so the
first-order deadweight loss of a seller-remitted tax is only the ex-post
return-margin term.

Whether the omitted ex-post term adds to or offsets that market-quantity fall
hinges on how the return rate responds to the price change.  When
pass-through is positive and the return rate rises with price, more units are
returned and handling costs rise, so the classical formula underestimates the
welfare loss.  When the return rate instead falls with price, the ex-post
term is a gain that partially offsets the ex-ante loss, so the classical
formula overestimates it.  When pass-through is negative, that same omitted
ex-post gain is what delivers the Pareto improvement above.  An analyst who
ignores returns then not only mismeasures the magnitude but gets the wrong
sign.

To sign that bias, we provide a simple diagnostic, from the return-rate
elasticity and the ratio of handling to list price, for when the classical
no-return pass-through formula over- or underestimates the correct
returns-adjusted pass-through.  Parallel conditions in the return-rate,
supply, and demand elasticities sign when the classical welfare formula over-
or understates the true first-order welfare change, or reverses its sign.
Those welfare-improving outcomes stand in sharp contrast to the classical
no-return benchmark, where pass-through is always positive and such gains are
impossible.

%=============================================================================
\section{Setting}
%=============================================================================

A homogeneous product sells at list price $p$.  Consumers indexed by $i$ differ
in return hassle $\kappa_i$, the disutility of sending the good back, and in
the post-purchase match distribution $G_i$ over match quality $u$.  After
buying, type $i$ draws $u\sim G_i$ and keeps the item when
$u-p\geq -\kappa_i$, earning $u-p$, and returns it otherwise, earning
$-\kappa_i$.  These payoffs are written for the case in which the refund on a
return equals the purchase price; that coincidence is relaxed under buyer
remittance in Section~\ref{sec:tax-neutrality}.

Expected surplus from buying at price $p$ is therefore
\[
  cs_i(p)=\mathbb{E}_{u\sim G_i}\bigl[\max(u-p,\,-\kappa_i)\bigr].
\]
On the purchase margin, all of this consumer-side heterogeneity collapses into
willingness to pay $v_i$, defined by $cs_i(v_i)=0$: type $i$ buys if and only
if $p\leq v_i$.  Across types, $v_i$ has distribution $F$ with density $f$, so
gross demand is $D(p)=1-F(p)$.  Of the units sold at $p$, a share $r(p)$ are
returned.

How $r(p)$ responds to price is central for what follows, so we unpack its
derivative.  Type $i$'s return probability is
\[
  \phi_i(p)=\Pr(u<p-\kappa_i).
\]
Willingness to pay settles who buys, but not who returns: buyers with the same
$v_i$ may still differ in $\kappa_i$ and $G_i$.  Averaging within WTP cells
gives
\[
  \phi(v,p)\equiv\mathbb{E}\bigl[\phi_i(p)\bigm|v_i=v\bigr],
\]
which depends on $v$ whenever the joint distribution of $(\kappa_i,G_i)$ varies
across valuation levels.  The aggregate return rate is the mean of these cell
averages among those who buy at $p$,
\[
  r(p)=\frac{1}{D(p)}\int_p^\infty \phi(v,p)\,dF(v).
\]
Differentiating and using $D'(p)=-f(p)$ splits the response into two terms,
\[
  r'(p)
  =
  \underbrace{
    \frac{1}{D(p)}\int_p^\infty \frac{\partial \phi(v,p)}{\partial p}\,dF(v)
  }_{\text{Behavioral channel}}
  +
  \underbrace{
    \frac{f(p)}{D(p)}\bigl(r(p)-\phi(p,p)\bigr)
  }_{\text{Selection channel}},
\]
where $\phi(p,p)$ is the average return probability among types at the
purchase margin $v_i=p$.

The first term is behavioral.  Among consumers who remain in the market, a
higher $p$ raises the keep threshold and makes return strictly more likely.
The second is selection.  A higher price truncates the buyer pool at the
margin, so composition changes whenever marginal buyers differ in return
propensity from the average inframarginal buyer.  For arbitrary
$(\kappa_i,G_i)$ that selection term is unsigned.  Under the additive match
$u=v_i+\varepsilon$ with $\varepsilon$ and return hassle both independent of
$v_i$, however, $\phi(v,p)$ decreases in $v$.  Marginal buyers at $v_i=p$ then
return more often than inframarginal buyers, so $\phi(p,p)\geq r(p)$ and
selection is non-positive.  The two forces therefore pull opposite ways.
Which dominates is empirical, and we impose no sign restriction on $r'(p)$
below.

On the firm side the market is perfectly competitive.  Producer $j$ takes $p$
as given, chooses output at cost $C_j(q_j)$, and incurs a net return handling
cost $h$ on each returned unit---gross handling cost minus scrap value.  Net
handling may be negative, but we maintain $p+h>0$, which holds whenever scrap
value does not exceed list price.

%=============================================================================
\section{Consumer side}
%=============================================================================

We turn first to consumer surplus and show that only the aggregates $D$ and
$r$ enter.  Write $cs(v,p)$ for expected surplus at list price $p$ among types
with willingness to pay $v$, so $cs(v,v)=0$.  Consumer surplus integrates surplus
over all types who buy,
\begin{equation}
  CS(p)=\int_p^\infty cs(v,p)\,dF(v).
\end{equation}
Differentiating in $p$ gives an extensive-margin term from types entering or
leaving the purchase margin and an intensive-margin term from lower surplus
among inframarginal buyers,
\[
  CS'(p)
  = -cs(p,p)\,f(p)+\int_p^\infty cs_p(v,p)\,dF(v),
\]
where $cs_p(v,p)\equiv \partial cs(v,p)/\partial p$.  On the extensive margin,
$cs(p,p)=0$, so entry and exit contribute nothing at first order.  On the intensive margin, a higher price harms a buyer only when she keeps, so the cell average is minus the keep probability in that WTP cell,
\[
  cs_p(v,p)
  =
  -\Pr(\text{keep}\mid v_i=v).
\]
Aggregating over buyers who purchase then gives
\[
  CS'(p)=\int_p^\infty cs_p(v,p)\,dF(v)
  = -D(p)\bigl(1-r(p)\bigr),
\]
by the definitions of $D(p)$ and $r(p)$.  A small price change therefore satisfies
\begin{equation}
  dCS = -D(p)\bigl(1-r(p)\bigr)\,dp.
  \label{eq:dCS}
\end{equation}

The extensive- and intensive-margin decomposition makes the local formula
transparent.  Consumer surplus changes depend only on the aggregates $D(p)$ and
$r(p)$, both at each price and integrated along any path,
\[
  \Delta CS = -\int_{p_0}^{p_1} D(p)\bigl(1-r(p)\bigr)\,dp.
\]
Micro heterogeneity in price responses can be rich.  Price $p$ enters each
$cs_i(p)$ nonlinearly and can shift participation, return decisions, and
inframarginal surplus differently across types.  Yet once $(D,r)$ are given, this
microstructure is irrelevant for $dCS$ and $\Delta CS$.  Two economies with
the same paths of $D(p)$ and $r(p)$ deliver the same consumer surplus loss,
whether high-$v_i$ types adjust return behavior more steeply than low-$v_i$
types or not.  Only kept units enter at first order, through $1-r(p)$, not
gross purchases $D(p)$ alone.

Let $\bar{p}$ denote the highest willingness to pay in the population, so
$D(\bar{p})=0$ and $CS(\bar{p})=0$.  Integrating the local formula then
rewrites consumer surplus as area under kept volume,
\begin{equation}
  CS(p)=\int_p^{\bar{p}} D(z)\bigl(1-r(z)\bigr)\,dz.
\end{equation}

%=============================================================================
\section{Supply side}
%=============================================================================

On the firm side, returns make list price an imperfect summary of producer
payoffs.  A gross sale brings revenue $p$ if the consumer keeps the item and
costs refund $p$ plus net handling $h$ if the item is returned.  Among gross
sales at list price $p$, share $r(p)$ are returned, so expected revenue per
gross unit---the \emph{effective margin}---is
\begin{equation}
  \mu(p,h) \equiv p - r(p)(p+h).
  \label{eq:mu-def}
\end{equation}
Under perfect competition, firm $j$ takes $p$ as given and therefore faces a
fixed number $\mu$ in its supply decision.  At equilibrium, $\mu=\mu(p,h)$.

Given $\mu$, firm $j$ chooses output to solve
\[
  \max_{q_j \geq 0}\;\; \mu\,q_j - C_j(q_j),
  \qquad
  \text{FOC: } C_j'(q_j)=\mu \;\;(\text{if } q_j>0),
\]
with $C_j(0)=0$.  Write $q_j(\mu)$ for firm $j$'s supply, set to zero when
$\mu$ is below its minimum marginal cost.  We assume each $C_j$ has strictly
increasing marginal cost, so $q_j'(\mu)>0$ for active firms.

Summing across firms defines market supply and industry profit,
\begin{equation}
  S(\mu) \equiv \sum_j q_j(\mu),
  \qquad
  PS(\mu) \equiv \sum_j \bigl[\mu\,q_j(\mu)-C_j\bigl(q_j(\mu)\bigr)\bigr].
  \label{eq:agg-def}
\end{equation}
Under that assumption, $S'(\mu)>0$.  Market supply therefore rises with the
effective margin, but not necessarily with list price.  Firms care about
$\mu$, and $\mu(p,h)$ can fail to rise with $p$ when refunds and handling
eat into the list-price gain.  Which firms are active at each $\mu$
determines both the market quantity and total profit.

Differentiating the profit of firm $j$ with respect to $\mu$ gives
\[
  \frac{d}{d\mu}\Bigl[\mu\,q_j-C_j(q_j)\Bigr]
  = q_j + q_j'(\mu)\,\bigl[\mu-C_j'(q_j)\bigr].
\]
Active firms satisfy $C_j'(q_j)=\mu$ at their chosen output, so the adjustment
term $q_j'(\mu)\,[\mu-C_j'(q_j)]$ is zero---equivalently, the envelope theorem
applied to firm $j$'s profit problem gives $d\pi_j/d\mu=q_j(\mu)$.  Summing
over firms,
\[
  PS'(\mu)=\sum_j q_j(\mu)=S(\mu).
\]
With no production at $\mu=0$, it follows that industry profit equals the area
to the left of the market supply curve,
\begin{equation}
  PS(\mu)=\int_0^\mu S(z)\,dz.
  \label{eq:PS-integral}
\end{equation}
At equilibrium $\mu=\mu(p,h)$,
$PS=\int_0^{\mu(p,h)}S(z)\,dz$.

%=============================================================================
\section{Failure of tax neutrality}
%=============================================================================
\label{sec:tax-neutrality}

With firms responding to the effective margin rather than the list price, a
natural question is whether the classical result---that it does not matter
which party remits a unit tax---survives.  \citet{weyl2013pass}
establish this neutrality without returns.  Write $p$ for the list
price and $\tau$ for a per-unit tax.
Write superscripts $s$ and $b$ for seller and buyer remittance respectively.
Under seller remittance the seller refunds the tax-inclusive list price
and recovers none of $\tau$ from the government, so purchase price and return
recovery both equal list price $p$.  The return rate is therefore the earlier
$r(p)$, while the seller's effective margin subtracts the unreclaimed tax,
\[
  \mu^{s}(p,\tau,h)=p-\tau-r(p)(p+h)=\mu(p,h)-\tau,
\]
and clearing is
\begin{equation}
  D(p) = S\bigl(\mu^{s}(p,\tau,h)\bigr).
  \label{eq:clear-seller}
\end{equation}
Consumer surplus depends on list price $p$ alone, as in Section~3, and producer
surplus is the area under supply at $\mu^{s}$,
\[
  CS^{s}=CS(p),
  \qquad
  PS^{s}=\int_0^{\mu^{s}}S(z)\,dz.
\]
A tax therefore reaches consumers only by moving the equilibrium list
price.  Differentiating at $\tau=0$, where $\mu^{s}_{\tau}=-1$ at fixed $p$ and
$\rho^{s}_{\tau}\equiv dp^{s}/d\tau$,
\begin{equation}
  \frac{dCS^{s}}{d\tau}
  =-Q\bigl(1-r\bigr)\rho^{s}_{\tau},
  \qquad
  \frac{dPS^{s}}{d\tau}
  =Q\bigl(\mu_p\rho^{s}_{\tau}-1\bigr).
  \label{eq:dCS-dPS-seller}
\end{equation}

\subsection{Buyer remittance: purchase price and recovery}

Buyer remittance breaks that coincidence.  The buyer pays $p$ in
total, of which $\tau$ goes to the government and $p-\tau$ to the seller.  On
a return the seller refunds that same pre-tax amount $p-\tau$, and the buyer
recovers none of $\tau$ from the government.  Purchase is still at list price
$p$, but the buyer's recovery is only
\[
  R=p-\tau.
\]
Expected surplus from buying therefore depends on the tax through the return
option,
\[
  cs_i(p,\tau)
  =\mathbb{E}_{u\sim G_i}\bigl[\max(u-p,\,-\tau-\kappa_i)\bigr].
\]
For each fixed $\tau$, type $i$'s willingness to pay $v_i(\tau)$ is defined by
$cs_i\bigl(v_i(\tau),\tau\bigr)=0$: she buys if and only if $p\leq v_i(\tau)$.
Differentiating the identity $cs_i\bigl(v_i(\tau),\tau\bigr)=0$ gives
\[
  v_i'(\tau)
  =-\frac{\partial cs_i/\partial\tau}{\partial cs_i/\partial p}.
\]
At list price $p$, $\partial cs_i/\partial p=-\Pr(\text{keep})$ and
$\partial cs_i/\partial\tau=-\Pr(\text{return})$, so
\[
  v_i'(\tau)
  =-\frac{\Pr(\text{return})}{\Pr(\text{keep})}
  \leq 0
\]
when evaluated at $p=v_i(\tau)$ (and $v_i'(\tau)=0$ if type $i$ never
returns).  A higher tax weakly lowers every type's willingness to pay by
worsening the return option.

Across types, $v_i(\tau)$ has distribution $F(\,\cdot\,;\tau)$.  At $\tau=0$,
recovery equals the list price, $cs_i(p,0)$ coincides with the earlier $cs_i(p)$, and
$v_i(0)$ recovers the earlier $v_i$, so $F(\,\cdot\,;0)=F$ with density $f$.
Because $v_i'(\tau)\leq 0$ for every $i$, a rise in $\tau$ shifts the law of
willingness to pay down in the first-order stochastic dominance (FOSD) sense:
for $\tau'>\tau$ and every $p$,
\[
  F(p;\tau')\geq F(p;\tau).
\]
Locally, writing $f(\,\cdot\,;\tau)$ for the density of $F(\,\cdot\,;\tau)$, the
CDF rises by the flux of types across $p$.  Substituting
$v_i'(\tau)=-\Pr(\text{return})/\Pr(\text{keep})$ at the purchase margin
$v_i(\tau)=p$,
\[
  \frac{\partial F(p;\tau)}{\partial\tau}
  =\mathbb{E}\!\left[
    \frac{\Pr_i(\text{return})}{\Pr_i(\text{keep})}
    \,\middle|\,v_i(\tau)=p
  \right]f(p;\tau)
  \geq 0.
\]
Gross demand is the upper tail of this distribution,
\[
  D(p,\tau)=1-F(p;\tau),
\]
so $D(p,0)=D(p)$ and, at fixed $p$,
\[
  D_\tau
  \equiv\frac{\partial D}{\partial\tau}
  =-\frac{\partial F(p;\tau)}{\partial\tau}
  =-\mathbb{E}\!\left[
    \frac{\Pr_i(\text{return})}{\Pr_i(\text{keep})}
    \,\middle|\,v_i(\tau)=p
  \right]f(p;\tau)
  \leq 0.
\]

Type $i$ keeps if and only if $u\geq R-\kappa_i$.  Write $\tilde{r}(p,R)$ for
the return rate among buyers when the list price is $p$ and recovery is $R$, with
the buyer set allowed to depend on both.  Under buyer remittance,
\begin{equation}
  \tilde{r}(p,p-\tau)
  =\frac{1}{D(p,\tau)}\int_p^\infty \phi(v,p-\tau)\,dF(v;\tau).
\end{equation}
The earlier return rate is the special case $r(p)=\tilde{r}(p,p)$.  Write
\[
  r_R\equiv \tilde{r}_R(p,p)
\]
for the partial of $\tilde{r}$ with respect to recovery $R$, holding the
list price $p$ fixed, evaluated at $R=p$.  The total derivative $r'(p)$
defined earlier raises purchase price and recovery together along $R=p$, so
\[
  r'(p)=\tilde{r}_p(p,p)+r_R.
\]
A buyer tax shifts $R$ but not the list price, so the relevant return
response is $r_R$, not $r'(p)$.

On the firm side, producer surplus is still the area under supply, but
evaluated at the buyer-remittance effective margin
\begin{equation}
  \tilde{\mu}^{b}(p,\tau,h)
  =(p-\tau)-\tilde{r}(p,p-\tau)\,(p-\tau+h),
  \label{eq:mu-tilde}
\end{equation}
rather than at $\mu(p,h)$.  Clearing and producer surplus are therefore
\begin{equation}
  D(p,\tau)=S\bigl(\tilde{\mu}^{b}(p,\tau,h)\bigr),
  \qquad
  PS=PS\bigl(\tilde{\mu}^{b}(p,\tau,h)\bigr)
  =\int_0^{\tilde{\mu}^{b}(p,\tau,h)}S(z)\,dz.
  \label{eq:clear-buyer}
\end{equation}
Differentiating at $\tau=0$, where $D(p,0)=D(p)$ and $Q\equiv D(p)$,
\begin{equation}
  \frac{dPS}{d\tau}
  =Q\bigl(\mu_p\rho^{b}_{\tau}+\tilde{\mu}^{b}_{\tau}\bigr),
  \label{eq:dPS-buyer}
\end{equation}
where $\rho^{b}_{\tau}\equiv dp^{b}/d\tau$ is the pass-through of the tax into
list price and
\begin{equation}
  \tilde{\mu}^{b}_{\tau}
  =-\bigl(1-r-(p+h)\,r_R\bigr)
\end{equation}
is the partial of $\tilde{\mu}^{b}$ with respect to $\tau$ at fixed
$p$.\footnote{At $\tau=0$, $\tilde{\mu}^{b}(p,0,h)=\mu(p,h)$.  Holding
  $\tau$ fixed, recovery $R=p-\tau$ rises one-for-one with the list price, so
  $\partial\tilde{\mu}^{b}/\partial p|_{\tau=0}=1-r-r'(p)(p+h)=\mu_p$.  The
  $\tau$-partial at fixed $p$ instead shifts only $R$, which is why it
  involves $r_R$ rather than $r'(p)$.}

Write $cs(v,p,\tau)$ for expected surplus among types with willingness to pay
$v_i(\tau)=v$, and set
\begin{equation}
  CS(p,\tau)
  =\int_p^\infty cs(v,p,\tau)\,dF(v;\tau).
\end{equation}
The total derivative with respect to $\tau$ is
\[
  \frac{dCS}{d\tau}
  =\frac{\partial CS}{\partial p}\Bigm|_{\tau}\rho^{b}_{\tau}
  +\frac{\partial CS}{\partial \tau}\Bigm|_{p}.
\]
For the price partial, differentiating the integral in $p$ yields an
extensive-margin term from types entering or leaving the purchase margin
and an intensive-margin term among remaining buyers.  As before, surplus is
zero at the margin, so the extensive-margin effect is zero at first order;
the intensive term harms keepers only.

For the tax partial at fixed list price, sum $cs_i(p,\tau)$ over buyers
($v_i(\tau)\ge p$).  Types who enter or leave the purchase margin have
$cs_i(p,\tau)=0$, so they contribute nothing at first order.  Among buyers,
$\partial cs_i/\partial\tau=-\Pr_i(\text{return})$: the tax lowers the return
payoff one-for-one.  Aggregating over buyers gives
$\partial CS/\partial\tau=-D(p,\tau)\,\tilde{r}(p,p-\tau)$.\footnote{Differentiating
  $\int cs\,dF(v;\tau)$ appears to add a term from $F(\,\cdot\,;\tau)$ shifting
  among buyers with $v>p$.  That is only a relabeling of types.  Among buyers
  the tax changes surplus only through return probabilities, already averaged
  into $-D\tilde{r}$.}  At
$\tau=0$,\begin{equation}
  \frac{\partial CS}{\partial p}\Bigm|_{\tau}
  =-D(p)\bigl(1-r(p)\bigr),
  \qquad
  \frac{\partial CS}{\partial \tau}\Bigm|_{p}
  =-D(p)\,r(p),
\end{equation}
and therefore
\begin{equation}
  \frac{dCS}{d\tau}
  =-Q\bigl(1-r\bigr)\rho^{b}_{\tau}-Qr.
  \label{eq:dCS-buyer}
\end{equation}
Compared with \eqref{eq:dCS-dPS-seller}, the producer formula differs only in
the direct tax effect on the effective margin:
$\tilde{\mu}^{b}_{\tau}$ replaces $-1$.  The consumer formula adds the
direct returner term $-Qr$, because a tax at fixed $p$ lowers recovery
one-for-one and therefore harms those who return.

\subsection{Neutrality fails}

The two regimes clear through $D(p)=S(\mu^{s})$ and
$D(p,\tau)=S(\tilde{\mu}^{b})$.  Remittance is therefore neutral only if
demand and the effective margins match at a common list price,
$D(p)=D(p,\tau)$ and $\mu^{s}(p,\tau,h)=\tilde{\mu}^{b}(p,\tau,h)$.
Canceling $p-\tau$ in the margins,
\begin{result}[Failure of tax neutrality with returns]
  Seller and buyer taxes define the same equilibrium if and only if
  $D(p)=D(p,\tau)$ and
  $r(p)(p+h)=\tilde{r}(p,p-\tau)\,(p-\tau+h)$.
\end{result}
Demand matches locally if and only if $D_{\tau}=0$.  That partial is in
general nonzero unless types at the purchase margin $v_i(\tau)=p$ never
return.  If the marginal buyer returns, $D_{\tau}<0$, so
$D(p,\tau)<D(p)$ for a small positive tax, the demand sides of the two
clearing equations differ, and neutrality fails.

%=============================================================================
\section{Pass-through}
%=============================================================================

Equilibrium satisfies $D(p)=S(\mu(p,h))$.  Write $Q\equiv D(p)$.  We
characterize pass-through into list price for three shocks: a uniform
one-dollar parallel shift in every firm's marginal production cost (indexed by
$c$), a one-dollar increase in handling cost $h$, and a unit tax under seller
and buyer remittance.  Production cost shifts market supply at a given
effective margin $\mu$, while handling lowers $\mu(p,h)$ at a fixed list
price $p$.

For production and handling, write the clearing condition as
\begin{equation}
  m(p,x)\equiv D(p)-S(\mu(p,x),x)=0,
  \qquad x\in\{c,h\}.
\end{equation}
Differentiating while keeping $m=0$ defines pass-through $\rho_x\equiv dp/dx$ by
\begin{equation}
  \rho_x=-\frac{m_x}{m_p},
  \label{eq:pass-through-general}
\end{equation}
and expanding $m=D-S$ gives
\begin{align}
  m_p &= D'(p)-S'(\mu)\,\mu_p,
  &
  m_x &= -S'(\mu)\,\mu_x-S_x,
\end{align}
hence
\begin{equation}
  \rho_x
  = \frac{-S'(\mu)\,\mu_x-S_x}{S'(\mu)\,\mu_p-D'(p)}.
  \label{eq:pass-through}
\end{equation}
The two shocks differ in which margin moves at fixed $p$ or fixed $\mu$,
\begin{align}
  \mu_c &= 0,
  &
  S_c &= -S'(\mu),
  &
  \mu_h &= -r(p),
  &
  S_h &= 0.
\end{align}
The production term is $S_c=-S'(\mu)$ because a one-dollar parallel rise in
every firm's marginal cost is equivalent, at fixed $\mu$, to evaluating market
supply at $\mu-1$.  Handling moves $\mu$ by only $r(p)$ per dollar relative to
that unit production shift, so its numerator in \eqref{eq:pass-through} is the
fraction $r(p)$ of the production numerator.

Turning to the unit tax, failure of tax neutrality requires separate list
prices $p^s$ and $p^b$ under seller and buyer remittance.  Write
$\rho^s_\tau\equiv dp^s/d\tau$ and $\rho^b_\tau\equiv dp^b/d\tau$.  Seller
remittance keeps demand at $D(p)$, so at $\tau=0$ \eqref{eq:pass-through}
with $S_\tau=0$ gives
\begin{equation}
  \rho^s_\tau
  = \frac{-S'(\mu)\,\mu^s_\tau}{S'(\mu)\,\mu_p-D'(p^s)}.
\end{equation}
Buyer remittance clears through $D(p,\tau)=S(\tilde{\mu}^{b})$.  Differentiating
at $\tau=0$, with $D_p\equiv\partial D/\partial p$ and
$D_\tau\equiv\partial D/\partial\tau$ evaluated at $\tau=0$,
\begin{equation}
  \rho^{b}_{\tau}
  =\frac{-S'(\mu)\,\tilde{\mu}^{b}_{\tau}+D_\tau}
  {S'(\mu)\,\mu_p-D_p}.
\end{equation}
At $\tau=0$ one has $D_p=D'(p)$ and $\mu^{s}_{\tau}=-1$, so
$\rho^{s}_{\tau}=\rho_c$.  The buyer rate is
\[
  \rho^{b}_{\tau}
  =-\tilde{\mu}^{b}_{\tau}\,\rho_c
  +\frac{D_\tau}{S'(\mu)\,\mu_p-D'(p)}.
\]
When demand does not respond directly to the tax ($D_\tau=0$), this collapses
to $\rho^{b}_{\tau}=-\tilde{\mu}^{b}_{\tau}\,\rho_c$.  When the purchase-price
selection channel in $r'(p)$ vanishes, $\tilde{\mu}^{b}_{\tau}=-\mu_p$ and the
older identification of a buyer tax with $-p$ returns (again if $D_\tau=0$).
If $\tilde{\mu}^{b}_{\tau}>0$ and $D_\tau=0$, buyer- and seller-remitted
pass-through carry opposite signs.  When returns are absent ($r=0$ and
$r_R=0$) and $D_\tau=0$, $\tilde{\mu}^{b}_{\tau}=-1$ and the two rates
coincide, recovering the classical neutrality benchmark.
\begin{result}[Pass-through into list price and remittance-dependent sign]
\label{res:pass-through}
  \textbf{(i)} Production-cost pass-through is
  \begin{equation}
    \rho_c
    = \frac{S'(\mu)}{S'(\mu)\,\mu_p-D'(p)}.
    \label{eq:rho-c}
  \end{equation}
  The remaining shocks connect to this object by
  \begin{align}
    \rho_h &= r(p)\,\rho_c,
    &
    \rho^{s}_{\tau} &= \rho_c,
    &
    \rho^{b}_{\tau}
    &= -\tilde{\mu}^{b}_{\tau}\,\rho_c
    +\frac{D_\tau}{S'(\mu)\,\mu_p-D'(p)}.
    \label{eq:rho-links}
  \end{align}
  \textbf{(ii)} If $D_\tau=0$ and $\tilde{\mu}^{b}_{\tau}>0$, then
  \[
  \operatorname{sgn}(\rho^{b}_{\tau})=-\operatorname{sgn}(\rho_c)
  \]
  provided $\rho_c\neq 0$.  A unit tax therefore moves the equilibrium
  list price in opposite directions depending on which side of the
  market remits it.
\end{result}
The condition $\tilde{\mu}^{b}_{\tau}>0$ requires the recovery response of
returns to be strong enough that $\tilde{\mu}^{b}$ rises when $\tau$ rises at
fixed $p$: $\tilde{\mu}^{b}_{\tau}>0$ exactly when $r_R(p+h)>1-r$.  When
also $D_\tau=0$, buyer remittance raises the effective margin while seller
remittance lowers it, shifting supply outward and inward respectively.  The
list price must therefore adjust in opposite directions to restore
equilibrium.  Which regime raises the price and which lowers it depends on
the sign of $\rho_c$, which in turn hinges on the relative magnitudes of the
supply and demand elasticities---a point developed in the next subsection.
\subsection{Elasticity representation}

The same pass-through formulas can be rewritten in elasticities, which makes
the role of $\mu_p$ transparent.  Write $\mu_p\equiv d\mu/dp$ for the response
of the effective margin to list price, and define demand and supply elasticities
with respect to list price as usual by
\[
\varepsilon_D(p)\equiv\frac{p D'(p)}{D(p)}<0,
\qquad
\varepsilon_S(p)\equiv\frac{dS(\mu(p))}{dp}\frac{p}{S(\mu(p))}
=S'(\mu)\,\mu_p\,\frac{p}{S(\mu)}.
\]
Here $\varepsilon_S(p)$ is the observed elasticity of market supply with
respect to the clearing price.  Without returns, $\mu_p=1$ and
$\varepsilon_S>0$.  With returns, $\mu_p$ can become negative, in which case
$\varepsilon_S<0$ and a higher list price reduces the quantity firms are
willing to supply.

Without returns, supply is a function of list price, so $S'=dS/dp$.  A
one-dollar cost increase cuts supply by $S'$ units at the original price,
creating an excess-demand gap of that size.  Raising list price by one dollar
closes the gap from both sides.  Demand falls by $|D'|$ and supply recovers by
$S'$.  Closing a gap of size $S'$ therefore requires a price rise of
$S'/(S'+|D'|)$, so standard pass-through equals
$\varepsilon_S/(-\varepsilon_D+\varepsilon_S)$.

With returns, supply is a function of the effective margin, so
$S'(\mu)=dS/d\mu$ rather than $dS/dp$.  A one-dollar production-cost increase
is equivalent to lowering $\mu$ by one dollar, so at fixed list price supply
falls by $S'(\mu)$ units, again creating an excess-demand gap of that size.
Raising list price by one dollar closes the gap at rate
$S'(\mu)\,\mu_p-D'(p)$, so
\[
\rho_c=\frac{S'(\mu)}{S'(\mu)\,\mu_p-D'(p)}.
\]
The cost shock lowers $\mu$ by one dollar, while a dollar of list price raises
$\mu$ by $\mu_p$.  A higher $\mu_p$ means a one-dollar price increase restores
more of the effective margin---and therefore more quantity supplied---than the
one-dollar cost shock removes, so a smaller list-price movement closes the
excess-demand gap.  Measuring supply against list price through
$\varepsilon_S$ puts the factor $1/\mu_p$ in the numerator,
\begin{equation}
  \rho_c
  =
  \frac{\varepsilon_S(p)/\mu_p}
       {-\varepsilon_D(p)+\varepsilon_S(p)}.
  \label{eq:rho_c_elasticity_final}
\end{equation}
When $\mu_p<0$, the supply curve expressed as a function of the list price
$p$ is locally decreasing.  In that region the sign of $\rho_c$ is no longer
guaranteed positive.  It depends on the relative magnitudes of $\varepsilon_S$
and $\varepsilon_D$.  This opens the door to the counterintuitive outcome that
an increase in marginal cost reduces the equilibrium price.
\begin{result}[Negative pass-through]
\label{res:negative-pt}
  We have $\rho_c<0$ if and only if $\varepsilon_S<0$ ($\mu_p<0$) and
  $|\varepsilon_S|>-\varepsilon_D$.
\end{result}
In words, negative pass-through occurs when supply slopes downward in list
price---the behavioral return response dominates---and supply is steeper than
demand.  Under these conditions, a cost increase creates excess demand that
can only be resolved by lowering the price, an outcome impossible in the
classical model.
\begin{proof}
By (\ref{eq:rho_c_elasticity_final}),
\[
  \rho_c = \frac{\varepsilon_S/\mu_p}{-\varepsilon_D+\varepsilon_S}.
\]
Since
$\varepsilon_S=S'(\mu)\mu_p\,p/S(\mu)$, the terms $\varepsilon_S$ and $\mu_p$
have the same sign, so $\varepsilon_S/\mu_p>0$.  Hence $\rho_c<0$ if and only
if
\[
  -\varepsilon_D+\varepsilon_S<0,
\]
which is equivalent to $\varepsilon_S<0$ and
$|\varepsilon_S|>-\varepsilon_D$.  Because $\varepsilon_S$ and $\mu_p$ have the
same sign, $\varepsilon_S<0$ is equivalent to $\mu_p<0$.
\end{proof}
The logic is as follows.  With $\varepsilon_S<0$ ($\mu_p<0$), a higher list price lowers the
quantity firms supply, so a leftward shift of supply---caused by a rise in
$c$---creates excess demand at the initial $p$.  Raising the price to eliminate
the excess would further contract supply and worsen the imbalance; instead,
the price must fall.  A decline in $p$ increases both the quantity demanded
and, because supply slopes downward, the quantity supplied.  When
$|\varepsilon_S|>-\varepsilon_D$, the supply curve is steeper than the demand
curve, meaning that the supply response to a price cut is larger than the
demand response.  The price decline therefore eliminates the excess demand and
leads to a new equilibrium at a lower list price.  This sign reversal cannot
occur in the classical model, where $\mu_p=1$ and supply is everywhere
upward-sloping, highlighting
a qualitative failure of the no-return benchmark.

Even when $\mu_p>0$ and pass-through stays positive, the returns-adjusted
formula differs from the no-return case whenever $\mu_p\neq 1$.  Comparing the
two signs the bias for an analyst who ignores returns.
Whether $\mu_p$ exceeds one depends on how steeply the return rate falls with
price.  Differentiating $\mu=p-r(p)(p+h)$ gives
\[
\mu_p=1-r(p)-r'(p)(p+h).
\]
Write $\varepsilon_r(p)\equiv p\,r'(p)/r(p)$ for the elasticity of the return rate
in list price.  Then $\mu_p>1$ if and only if $\varepsilon_r(p)<-p/(p+h)$.
That is, selection must outweigh the behavioral rise in individual return
probabilities by enough that aggregate $r$ falls steeply in $p$.  When instead
the behavioral response dominates, $\varepsilon_r$ is less negative (or
positive) and $\mu_p<1$.
\begin{result}
\label{res:pt-bias-er}
  Relative to the correct returns-adjusted $\rho_c$, the classical no-return
  pass-through formula overestimates pass-through when
  $\varepsilon_r(p)<-p/(p+h)$ and underestimates it otherwise.
\end{result}
The practical value is that the bias can be signed without estimating the full
returns-adjusted model.  The return-rate elasticity and the ratio of handling
to list price are enough.  The economics behind the threshold is familiar from
the selection and behavioral forces already in $\varepsilon_r$.  When selection
dominates, a higher list price lowers the return rate, so the
refund-and-handling drag on the effective margin shrinks as price rises.
The effective margin can then rise by more than one dollar per dollar of
list price, and a smaller price movement closes the excess-demand gap than
the classical formula predicts.  Costly handling strengthens that channel and
makes overestimation more likely.  Otherwise, the effective margin rises by
less than one dollar per dollar of list price, so list price must move by more
than the classical formula predicts to close the gap, and the classical
no-return formula underestimates pass-through.

%=============================================================================
\section{Local incidence}
%=============================================================================

Having established how prices respond to different shocks, we turn to local
incidence: how each shock's pass-through $\rho_x$ splits the resulting surplus
change between consumers and producers.  Throughout, $Q\equiv D(p)$ and
$r\equiv r(p)$ denote the equilibrium gross quantity and the return rate.
For shocks that reach consumers only through the list price---production
cost, handling cost, and seller remittance---the consumer-side analysis in
Section~3 gives
\begin{equation}
  \frac{dCS}{dx} = -Q(1-r)\rho_x,
  \label{eq:CSx_incidence}
\end{equation}
where $\rho_x\equiv dp/dx$.  Buyer remittance adds the direct recovery channel
derived above, so
\begin{equation}
  \frac{dCS}{d\tau^{b}}
  =-Q(1-r)\rho^{b}_{\tau}-Qr
\end{equation}
as in \eqref{eq:dCS-buyer}.

Producer surplus is $PS(\mu,x)=\int_0^{\mu} S(z,x)\,dz$.  Differentiating with
respect to $x$ and applying the equilibrium condition $S(\mu,x)=Q$ gives
\[
  \frac{dPS}{dx}
  = Q\frac{d\mu}{dx} + \Sigma_x,
\]
where $\Sigma_x\equiv\int_0^\mu (\partial S(z,x)/\partial x)\,dz$ captures the
direct shift in the area under the supply curve.  Expanding the total
derivative of the effective margin via the chain rule,
$\frac{d\mu}{dx} = \mu_p \rho_x + \mu_x$, yields
\begin{equation}
  \frac{dPS}{dx} = Q(\mu_p \rho_x + \mu_x) + \Sigma_x.
  \label{eq:PSx_incidence}
\end{equation}
Under buyer remittance this is \eqref{eq:dPS-buyer}, with
$\mu_x=\tilde{\mu}^{b}_{\tau}$ and $\rho_x=\rho^{b}_{\tau}$ from
\eqref{eq:rho-links} at $\tau=0$.

The local incidence ratio is $I_x \equiv (dCS/dx)/(dPS/dx)$.  For the three
price-only shocks,
\[
  I_x = \frac{-Q(1-r)\rho_x}{Q(\mu_p \rho_x + \mu_x) + \Sigma_x},
\]
and the components in Table~\ref{tab:incidence_components} reduce to a common
expression in $\rho_c$ and $\mu_p$.  The buyer-tax row uses
\eqref{eq:dCS-buyer}--\eqref{eq:dPS-buyer} instead, so its ratio differs by the
direct $-Qr$ term.

\begin{table}[htbp]
\centering
\begin{tabular}{lccccccc}
\hline
Shock $x$ & $\rho_x$ & $\mu_x$ & $\Sigma_x$ & $dCS/dx$ & $dPS/dx$ & $I_x$ \\
\hline
Production cost $c$ & $\rho_c$ & $0$ & $-Q$ & $-Q(1-r)\rho_c$ & $Q(\mu_p\rho_c - 1)$ & $\displaystyle \frac{(1-r)\rho_c}{1-\mu_p\rho_c}$ \\
Handling cost $h$ & $r\rho_c$ & $-r$ & $0$ & $-Q(1-r)r\rho_c$ & $Qr(\mu_p\rho_c - 1)$ & $\displaystyle \frac{(1-r)\rho_c}{1-\mu_p\rho_c}$ \\
Seller tax $\tau^{s}$ & $\rho_c$ & $-1$ & $0$ & $-Q(1-r)\rho_c$ & $Q(\mu_p\rho_c - 1)$ & $\displaystyle \frac{(1-r)\rho_c}{1-\mu_p\rho_c}$ \\
Buyer tax $\tau^{b}$ & $\rho^{b}_{\tau}$ & $\tilde{\mu}^{b}_{\tau}$ & $0$ & $-Q(1-r)\rho^{b}_{\tau}-Qr$ & $Q(\mu_p\rho^{b}_{\tau}+\tilde{\mu}^{b}_{\tau})$ & $\displaystyle \frac{-(1-r)\rho^{b}_{\tau}-r}{\mu_p\rho^{b}_{\tau}+\tilde{\mu}^{b}_{\tau}}$ \\
\hline
\end{tabular}
\caption{Pass-through, direct effects, surplus changes, and incidence ratio for
each shock.  Buyer-tax pass-through $\rho^{b}_{\tau}$ is as in
\eqref{eq:rho-links}.}
\label{tab:incidence_components}
\end{table}

\begin{result}
\label{res:incidence-ratio}
For production cost, handling cost, and seller remittance, the local incidence
ratio is identical:
\begin{equation}
  I_x = \frac{(1-r)\rho_c}{1-\mu_p \rho_c}>0.
  \label{eq:incidence_result}
\end{equation}
Under buyer remittance,
\begin{equation}
  I_{\tau^{b}}
  =\frac{-(1-r)\rho^{b}_{\tau}-r}{\mu_p\rho^{b}_{\tau}+\tilde{\mu}^{b}_{\tau}},
\end{equation}
with $\rho^{b}_{\tau}$ as in \eqref{eq:rho-links}.  When $D_\tau=0$, so that
$\rho^{b}_{\tau}=-\tilde{\mu}^{b}_{\tau}\,\rho_c$, this becomes
\[
  I_{\tau^{b}}
  =\frac{(1-r)\tilde{\mu}^{b}_{\tau}\rho_c-r}{\tilde{\mu}^{b}_{\tau}\bigl(1-\mu_p\rho_c\bigr)},
\]
which coincides with \eqref{eq:incidence_result} only in special cases (for
example $r=0$, or $\tilde{\mu}^{b}_{\tau}=-1$ with no direct recovery loss).
\end{result}

Result~\ref{res:incidence-ratio} establishes that the incidence ratio is
identical across the three price-only shocks, but the algebra alone can obscure
the deeper economic mechanism.  The key to understanding this result is that
each of those shocks is a scaled perturbation of the same equilibrium
geometry.  Production cost and seller-tax shocks shift the effective margin
$\mu$ down by one unit.  Handling cost shifts it down by $r$ units.  Locally,
the equilibrium condition $D(p)=S(\mu)$ is linear in the shock's direct effect
on $\mu$.  The resulting supply shift is proportional to $\mu_x$, and the price
adjustment $\rho_x$ that restores market clearing inherits that same
proportionality.  A buyer tax replaces the unit fixed-price response
$-1$ with $\tilde{\mu}^{b}_{\tau}$, but it also taxes returners directly at
fixed $p$, so the consumer side is no longer a pure multiple of $\rho_x$
and the common ratio breaks.

Because both $S(\mu)$ and $PS(\mu)$ are locally linear in $\mu$, the entire
equilibrium response to a price-only shock scales with the shock's effect on
the effective margin.  The shock's scalar therefore multiplies both consumer
and producer surplus changes in the same proportion and cancels out in the
incidence ratio.  The shock determines how far the equilibrium moves.  How that
movement is split between consumers and producers depends on the local demand
and supply slopes and on local $\mu_p$.\footnote{Consumers care about the
list-price change $\rho_x$.
Producers recover only through the effective margin, which a list-price rise
moves by $\mu_p\rho_x$.  Thus $\mu_p$ governs how much of each dollar of $p$
feeds into producer surplus.}
Incidence for those shocks is therefore fixed by market structure, not by
which side of the market is hit first.  What remains is why that common ratio
is positive.

The denominator $1-\mu_p\rho_c$ in \eqref{eq:incidence_result} is the per-unit
loss in producer surplus.  A \$1 cost shock initially reduces producer surplus
by $Q$.  The equilibrium list-price rise $\rho_c$ increases the effective
margin per unit by $\mu_p\rho_c$.  Because $dPS/d\mu=Q$ by the envelope
theorem, this marginal gain translates into a total recovery of
$Q\mu_p\rho_c$ in producer surplus.  The net loss is therefore
$Q(1-\mu_p\rho_c)$.  The numerator $(1-r)\rho_c$ is the corresponding
per-unit consumer loss on kept units.  The incidence ratio compares these two
per-unit burdens.  A small denominator leaves consumers
bearing almost the entire shock, while a negative denominator, $\mu_p\rho_c>1$,
means the recovery exceeds the cost, so producers are net gainers.  This
recovery factor can be written as
$\mu_p\rho_c=\varepsilon^S/(\varepsilon^S-\varepsilon^D)$.  This expression
closely mirrors the pass-through formula, and now we argue that
$\mu_p\rho_c>1$ if and only if $\rho_c<0$.  The condition
$\mu_p\rho_c>1$ requires $\varepsilon^S<0$, which implies $\mu_p<0$, and hence
$\rho_c<0$.  Conversely, $\rho_c<0$ requires $\varepsilon^S<0$ and
$|\varepsilon^S|>-\varepsilon^D$, which gives $\mu_p\rho_c>1$.  Therefore, the
condition governing producer gain from an increase in production cost is
exactly the same as that governing negative pass-through.

At first glance it may seem puzzling that, under negative pass-through, the
effective-margin recovery more than offsets the initial producer-surplus
loss of $Q$.  The price fall must do double duty to clear the market.  A one-dollar cost increase shifts the supply curve upward, creating
an initial excess demand gap.  To fill this gap, quantity supplied must
increase.  Since quantity supplied increases with the effective margin
$\mu$, restoring enough supply to cover the original gap already requires
$\mu$ to rise by one dollar.  But the price fall also increases quantity
demanded, creating additional excess demand on top of the initial gap.  To
satisfy both the original supply shortfall and this new demand, supply must
increase by more than the original shift, so the effective margin must rise
by more than one dollar.  Thus, $\mu_p\rho_c>1$ is precisely the
market-clearing requirement when pass-through is negative.

The preceding observation implies a striking result: the uniformly positive
incidence ratio means consumer and producer surplus always move together,
regardless of the sign of pass-through.  The direction of that common
movement nevertheless depends on both the shock type and the sign of
pass-through.  This yields a remarkable implication: a cost increase can be
Pareto-improving.  When $\rho_c<0$, raising production costs, handling costs,
or seller-remitted taxes lowers the list price, reduces return rates, and
saves enough handling costs to more than offset the direct resource drain;
both consumers and producers benefit.  Under the same condition,
buyer-remitted taxes reduce both surpluses, so the two tax regimes yield
opposite welfare outcomes.  This is a sharp violation of classical tax
neutrality.
When $\rho_c>0$, all four shocks reduce both surpluses.  Such gains are
impossible in the classical model, where pass-through is always positive,
underscoring how profoundly returns alter standard welfare analysis.

\subsection{First-order private surplus and deadweight loss}
\label{sec:DWL}

Local incidence separates the burden between consumers and producers.  Their
joint payoff---private surplus---is the sum of consumer and producer surplus,
$W\equiv CS+PS$.  Summing (\ref{eq:CSx_incidence}) and
(\ref{eq:PSx_incidence}) gives
\[
  \frac{dW}{dx}
  = -Q(1-r)\rho_x
    + Q(\mu_p\rho_x + \mu_x)
    + \Sigma_x
  = Q\bigl[\mu_x + \bigl(\mu_p-(1-r)\bigr)\rho_x\bigr]
    + \Sigma_x.
\]
Substituting $\mu_p=1-r-r'(p)(p+h)$ directly yields
\begin{equation}
  \frac{dW}{dx}
  = Q\bigl[\mu_x - r'(p)(p+h)\rho_x\bigr] + \Sigma_x,
  \label{eq:DWL_incidence}
\end{equation}
with $\mu_x$ and $\Sigma_x$ as in Table~\ref{tab:incidence_components}.
The term $-Q\,r'(p)(p+h)\rho_x$ is new relative to the classical model.  It
captures the efficiency consequence of the price-induced change in the return
rate.  Each returned unit consumes real handling resources $h$, so a change in
the return rate has a first-order effect on private surplus that is absent
when returns are shut down.

Without returns, $r=0$ and $r'=0$, so \eqref{eq:DWL_incidence} reduces to
\[
  \frac{dW}{dx} = Q\,\mu_x + \Sigma_x.
\]
For a unit production-cost increase, the supply curve shifts vertically upward
by one dollar at every quantity.  The surplus loss is therefore the area of
that shift---a rectangle of height one and width $Q$.  Algebraically, private
surplus is $W=CS(p)+PS(p-c)$.  Differentiating, the price-change terms
$-Q\rho_c$ and $+Q\rho_c$ cancel, leaving $dW/dc=-Q$: a first-order real
resource loss of $Q$ dollars.

For a unit tax the same cancellation gives $dW/d\tau=-Q$, but government
revenue $G=\tau Q(\tau)$ is transferred rather than lost.  Total surplus
$V=W+G$ therefore satisfies
$dV/d\tau=dW/d\tau+Q(\tau)+\tau Q'(\tau)=\tau Q'(\tau)$, which is zero at
$\tau=0$.  Locally, a cost shock has a first-order private-surplus effect of
$-Q$, while an infinitesimal tax has no first-order social loss---the dollar
is lost in one case and merely transferred in the other.

Product returns overturn that clean private-versus-transfer contrast.  The
return-rate term in \eqref{eq:DWL_incidence} is generally nonzero for every
shock, so production-cost increases, handling-cost increases, and unit taxes
under both remittance regimes all change private surplus at first order in a
way that is not a pure transfer between buyers and sellers.  We first track
these private-surplus changes for all four shocks.  For taxes, treasury
revenue remains a transfer; we add $G=\tau Q(\tau)$ back when measuring
deadweight loss below.

Plugging the values of $\mu_x$ and $\Sigma_x$ from
Table~\ref{tab:incidence_components} into
$dW/dx=Q\bigl[\mu_x-r'(p)(p+h)\rho_x\bigr]+\Sigma_x$, and using the
pass-through links from Result~\ref{res:pass-through}, yields
\begin{align}
  \frac{dW}{dc}
  &= -Q\,\Lambda,
  &
  \frac{dW}{dh}
  &= -Q\,r\,\Lambda,
  &
  \frac{dW}{d\tau^s}
  &= -Q\,\Lambda,
  &
  \frac{dW}{d\tau^b}
  &= -Q\,\mu_p\,\Lambda,
  \label{eq:DWL_shocks}
\end{align}
where the common factor is defined as
\begin{equation}
  \Lambda \equiv 1 + r'(p)(p+h)\,\rho_c.
  \label{eq:Lambda}
\end{equation}
Private surplus moves because consumers change either how much they buy or
whether they return what they bought.  The unit term in $\Lambda$ is the
mechanical effect on the purchase margin that remains even if the return rate
does not respond to price.  The term $r'(p)(p+h)\rho_c$ is the extra
efficiency cost on the keep-or-return margin after the price moves: each
change $r'(p)\rho_c$ in the return rate costs $p+h$ per unit sold.  We
therefore decompose the private-surplus changes in \eqref{eq:DWL_shocks} into
an ex-ante loss on the purchase margin and an ex-post loss on the
keep-or-return margin.  That split is simplest for a production-cost or
seller-tax shock, where $dW/dx=-Q\Lambda$; handling and buyer-remitted taxes
rescale both pieces by $r$ and $\mu_p$, respectively.

Freeze the aggregate return rate at its initial value $r_0\equiv r(p_0)$.
In the resulting ex-ante economy the effective margin is
\eqref{eq:mu-def} with $r(p)$ replaced by that constant,
\begin{equation}
  \mu^{\mathrm{EA}}(p)\equiv p-r_0(p+h).
  \label{eq:mu-EA}
\end{equation}
Applying \eqref{eq:DWL_incidence} then yields
$dW/dx\big|_{\mathrm{EA}}=-Q$, the classical rectangle from changing gross
quantity alone.  The ex-post change is that residual,
$dW/dx\big|_{\mathrm{EP}}\equiv dW/dx-dW/dx\big|_{\mathrm{EA}}$.  Since
$dW/dx=-Q\Lambda$ for these shocks,
\begin{equation}
  \frac{dW}{dx}\Big|_{\mathrm{EP}}
  =-Q\,r'(p)(p+h)\rho_c.
  \label{eq:DWL-EP}
\end{equation}
\begin{result}[Ex-ante and ex-post decomposition]
\label{res:EA-EP}
  For a unit production-cost or seller-tax shock,
  \begin{equation}
    \frac{dW}{dx}
    =-Q
    -Q\,r'(p)(p+h)\rho_c.
    \label{eq:DWL-decomp}
  \end{equation}
  The first term is the ex-ante loss on the purchase margin.  The second is
  the ex-post loss on the keep-or-return margin.  Each unit switched from keep
  to return costs $p+h$.  For a handling shock the two losses become $-Qr$ and
  $-Qr\,r'(p)(p+h)\rho_c$.  For a buyer-remitted tax they become $-Q\mu_p$ and
  $-Q\mu_p\,r'(p)(p+h)\rho_c$, so the ex-ante loss is not $-Q$ whenever
  $\mu_p\neq 1$.
\end{result}
Because consumer and producer surplus always move together, $\Lambda$ has the
same sign as $\rho_c$.  The no-return formula retains only the ex-ante
rectangle $-Q$, so the comparison with that benchmark is governed only by the
ex-post term $-Q\,r'(p)(p+h)\rho_c$.  That term may partially offset the
rectangle, amplify it, or reverse the private-surplus change.

The simplest case is when the selection channel in the return rate
dominates, so $r'(p)<0$.  A cost increase raises the list price (since
$r'<0$ forces $\mu_p>0$ and hence $\rho_c>0$), screens out frequent
returners, and lowers the return rate.  The ex-post term is then a gain that
partially offsets the ex-ante rectangle, so $0<\Lambda<1$.  It cannot reverse
the private-surplus change: $\rho_c>0$ keeps both consumer and producer
surplus falling, hence $\Lambda>0$.

When instead $r'(p)>0$, the bias hinges on the sign of pass-through.  If
$\rho_c>0$, the higher list price raises returns and adds handling costs, so
the ex-post loss amplifies the ex-ante rectangle and $\Lambda>1$.  That
amplification is larger when $r'(p)$ or $h$ is large.
% The same qualitative conclusion holds when $\mu_p<0$ but pass-through
% remains positive ($|\varepsilon_S|<-\varepsilon_D$): $\mu_p<0$ already
% requires a steep return response, so even a modest $\rho_c$ can push
% $\Lambda$ well above one.

If instead $\rho_c<0$, a production-cost or seller-tax increase lowers the
list price, so consumer surplus rises.  Positive incidence then raises
producer surplus as well---a Pareto improvement---and private surplus rises
by $Q|\Lambda|$ ($\Lambda<0$).  The gain is larger when
$r'(p)(p+h)\,|\rho_c|$ is large.  These three comparisons are summarized in
Result~\ref{res:welfare-bias}.
\begin{result}[No-return private-surplus bias]
\label{res:welfare-bias}
  Consider a unit production-cost or seller-tax increase.  Relative to the
  correct returns-adjusted change $dW/dc=-Q\Lambda$, the no-return formula
  $dW/dc=-Q$
  \begin{enumerate}
  \item[(i)] overestimates the private-surplus loss if and only if
  $r'(p)<0$ (which implies $\mu_p>0$ and hence $\rho_c>0$, so
  $0<\Lambda<1$);
  \item[(ii)] underestimates the private-surplus loss if and only if
  $r'(p)>0$ and $\rho_c>0$ (hence $\Lambda>1$);
  \item[(iii)] overestimates the private-surplus loss with the wrong
  sign---reporting a loss of $Q$ when the shock in fact raises private
  surplus by $Q|\Lambda|$---if and only if $r'(p)>0$ and $\rho_c<0$.
  \end{enumerate}
\end{result}
For production-cost and seller-tax shocks that bias is purely ex-post.  The
no-return formula counts the ex-ante rectangle correctly and errs only by
omitting the keep-or-return term.  Under buyer remittance both margins are
scaled by $\mu_p$,
\begin{equation}
  \frac{dW}{d\tau^b}
  =-Q\mu_p
  -Q\mu_p\,r'(p)(p+h)\rho_c
  =-Q\mu_p\Lambda,
\end{equation}
so relative to the returns-adjusted change $-Q\mu_p\Lambda$, the no-return
formula $-Q$ errs on both the ex-ante piece ($-Q$ vs $-Q\mu_p$) and the
ex-post piece.  When $\mu_p>0$, the no-return formula underestimates the
ex-ante loss if and only if $\mu_p>1$, and overestimates it otherwise.  That
cut is the same threshold as Result~\ref{res:pt-bias-er}: $\mu_p>1$ if and
only if $\varepsilon_r(p)<-p/(p+h)$.  Overestimation of pass-through by the
no-return formula is therefore underestimation of the ex-ante loss.

For the ex-post piece, as in Result~\ref{res:welfare-bias}(i), $r'(p)<0$
pins the term as a gain.  When selection dominates the behavioral response, the
higher list price from cost pass-through screens out consumers who return
frequently, so the return rate and the associated refund-and-handling cost
fall.  When $r'(p)>0$, the sign of the ex-post term depends on $\rho_c$ and
has no succinct formula.
\begin{corollary}[Buyer-remittance no-return bias]
\label{cor:welfare-bias-buyer}
  Relative to the correct returns-adjusted change $-Q\mu_p\Lambda$, the
  no-return prediction $dW/d\tau=-Q$ underestimates the private-surplus loss
  if and only if $\mu_p\Lambda>1$ and overestimates it if and only if
  $\mu_p\Lambda<1$.  The latter cut is equivalent to
  \begin{equation}
    \varepsilon_r(p)\,
    \frac{-\varepsilon_D}{\varepsilon_S-\varepsilon_D}
    >
    -\frac{p}{p+h}.
    \label{eq:V-buyer-cut}
  \end{equation}
\end{corollary}
Appendix~\ref{app:welfare-bias-primitives} proves the corollary.
Which sign of $\mu_p\Lambda-1$ obtains depends on the return-rate slope
and on pass-through.  Table~\ref{tab:buyer-bias-cases} summarizes the four
cases.  The first columns record the primitives and the two pass-throughs:
$\rho_c$ is cost pass-through, and buyer-tax pass-through into the list
price is $\rho^b_\tau=\mu_p\rho_c$.  The last three columns split the private
change $-Q\mu_p\Lambda$.  EA is the ex-ante (frozen-return) term $-Q\mu_p$: a
loss when $\mu_p>0$ and a gain when $\mu_p<0$.  EP is the ex-post keep-or-return
term $-Q\mu_p\,r'(p)(p+h)\rho_c$.  Bias vs $-Q$ reports the corollary cut:
overestimate when $\mu_p\Lambda<1$ and underestimate when $\mu_p\Lambda>1$.
Wrong sign is not a third category: it is the extreme case of overestimation
in which private surplus rises.  The no-return formula still reports a loss
of $Q$, while the true change is a gain, so $-Q$ both overestimates the loss
and gets its direction wrong.
\begin{table}[htbp]
\centering
\begin{tabular}{llccccc}
\hline
Case & Primitives & $\rho_c$ & $\rho^b_\tau$ & EA & EP & Bias vs $-Q$ \\
\hline
1 & $r'<0$ ($\Rightarrow\mu_p>0$) & $>0$ & $>0$ &
  loss & gain & over- or under-estimate \\
2 & $r'>0$, $\mu_p>0$ & $>0$ & $>0$ &
  loss & loss & overestimate \\
3 & $r'>0$, $\mu_p<0$, $|\varepsilon_S|<-\varepsilon_D$ & $>0$ & $<0$ &
  gain & gain & overestimate (wrong sign) \\
4 & $r'>0$, $\mu_p<0$, $|\varepsilon_S|>-\varepsilon_D$ & $<0$ & $>0$ &
  gain & loss & over- or under-estimate \\
\hline
\end{tabular}
\caption{Buyer-remittance bias relative to $-Q$.}
\label{tab:buyer-bias-cases}
\end{table}
\begin{itemize}
  \item Case~1 (selection).  Because $\mu_p>0$, a buyer-remitted tax lowers
  the effective margin at a fixed list price, so the ex-ante
  (frozen-return) private change $-Q\mu_p$ is a loss---larger than $-Q$ when
  $\mu_p>1$ and smaller when $\mu_p<1$.  The ex-post term is a gain as above.
  Overestimation ($\mu_p\Lambda<1$) further requires that
  $-\varepsilon_D/(\varepsilon_S-\varepsilon_D)$ not be too large---demand
  not too elastic relative to supply---so that pass-through remains large
  enough for the fall in the return rate to deliver a sufficient $p+h$
  saving.  When demand is very elastic, small positive pass-through weakens
  equilibrium screening and reduces this ex-post gain.
  \item Case~2 (mild behavioral).  Returns rise in price, but not enough to
  flip the effective-margin slope: $0<\mu_p<1-r$.  The ex-ante term is
  therefore a loss smaller than $Q$.  With $\rho_c>0$ we also have
  $\Lambda>1$, so the ex-post term is a loss as well.  The combined loss is
  nevertheless smaller than $Q$: writing
  $\mu_p\Lambda-1=-r-r'(p)(p+h)(1-\mu_p\rho_c)$,
  both summands are negative when $r'>0$ and $\rho_c>0$
  (hence $\mu_p\rho_c<1$).
  \item Case~3 (strong behavioral, positive cost pass-through).
  Returns rise in price strongly enough that $\mu_p<0$, so at a fixed
  list price a buyer-remitted tax \emph{raises} the seller's effective
  margin and the ex-ante term is a gain.  Demand is steeper than supply, so
  cost pass-through is positive and buyer-tax pass-through
  $\rho^b_\tau=\mu_p\rho_c$ is negative.  The list price therefore falls,
  and with $r'>0$ returns fall as well, making the ex-post term a gain.
  Private surplus rises ($\mu_p\Lambda<0$), while $-Q$ still reports a loss
  of $Q$---the extreme wrong-sign overestimate.
  \item Case~4 (strong behavioral, negative cost pass-through).
  Supply is steeper than demand, so cost pass-through is negative and
  buyer-tax pass-through $\rho^b_\tau=\mu_p\rho_c$ is positive.  The
  list price therefore rises, and with $r'>0$ returns rise as well,
  making the ex-post term a loss that opposes the ex-ante gain from
  $\mu_p<0$.  From the Case~2 identity
  $\mu_p\Lambda-1=-r-r'(p)(p+h)(1-\mu_p\rho_c)$,
  negative cost pass-through forces $\mu_p\rho_c>1$, so the second summand
  is positive and offsets $-r$.  The Case~4 primitives alone therefore do
  not pin whether $\mu_p\Lambda$ lies below or above one, nor whether
  private surplus falls or rises.
\end{itemize}

For tax shocks, social surplus adds back treasury revenue.  Because the
remitter recovers none of $\tau$ from the government on a return, revenue is
$G=\tau Q(\tau)$ under both remittance regimes, with $Q$ the equilibrium
gross quantity at that tax.  Write $V=W+G$.  Differentiating the product
gives $dG/d\tau=Q+\tau Q'$, so at $\tau=0$ we have $dV/d\tau=dW/d\tau+Q$.
Using \eqref{eq:DWL_shocks},
\begin{align}
  \frac{dV}{d\tau^s}
  &= -Q(\Lambda-1)
  = -Q\,r'(p)(p+h)\rho_c,
  &
  \frac{dV}{d\tau^b}
  &= Q\bigl(1-\mu_p\Lambda\bigr).
  \label{eq:DWL_V_tax}
\end{align}
Whether an infinitesimal tax creates a first-order social loss or gain
follows the same cuts as Result~\ref{res:welfare-bias} and
Corollary~\ref{cor:welfare-bias-buyer}.

%=============================================================================
\section{Global Incidence}
\label{sec:global-incidence}
%=============================================================================

Integrating the local derivatives $\frac{dCS}{dx}$ and $\frac{dPS}{dx}$ listed
in Table~\ref{tab:incidence_components} from $x_0$ to $x_1$ gives the global
surplus changes $\Delta CS$ and $\Delta PS$ for any shock
$x\in\{c,h,\tau^s,\tau^b\}$.  The global incidence ratio
$I_{x_0}^{x_1}\equiv \Delta CS/\Delta PS$ (when $\Delta PS\neq0$) follows from
the local relation $dCS/dx = I(x)\,dPS/dx$:
\begin{equation}
I_{x_0}^{x_1} = \int_{x_0}^{x_1} I(x)\,\omega(x)\,dx,
\label{eq:global-incidence-final}
\end{equation}
where $\omega(x)\equiv (dPS/dx)/\Delta PS$.  Although $I(x)>0$ at every point,
the global ratio need \emph{not} be positive: if $dPS/dx$ changes sign along
the path, the weights $\omega(x)$ take both signs, and the weighted average
can become negative.  Hence the classical sign-preserving property of local
incidence does not extend globally.

\clearpage
\appendix
\section{Welfare bias in primitives}
\label{app:welfare-bias-primitives}
Recall $\varepsilon_r(p)\equiv p\,r'(p)/r(p)$ and
$\mu_p=1-r(p)-r'(p)(p+h)$.  Then
$r'(p)(p+h)=\varepsilon_r(p)\,r(p)\,(p+h)/p$, so
\begin{equation}
  \mu_p>0
  \;\iff\;
  \varepsilon_r(p)
  <
  \frac{p}{p+h}\,\frac{1-r(p)}{r(p)},
  \label{eq:app-mu-er}
\end{equation}

Result~\ref{res:welfare-bias} becomes:
\begin{enumerate}
\item[(i)] overestimation iff $\varepsilon_r(p)<0$;
\item[(ii)] underestimation iff $\varepsilon_r(p)>0$ and either $\mu_p>0$,
or both $\mu_p<0$ and $|\varepsilon_S|\le -\varepsilon_D$;
\item[(iii)] overestimation with wrong sign (predicted loss, actual gain)
iff $\mu_p<0$ and $|\varepsilon_S|>-\varepsilon_D$.
\end{enumerate}
\begin{proof}
By Result~\ref{res:welfare-bias}, overestimation with $0<\Lambda<1$ is
$r'<0$, underestimation is $r'>0$ with $\rho_c>0$, and overestimation with
wrong sign is $r'>0$ with $\rho_c<0$.
Since $\varepsilon_r$ has the same sign as $r'$, write these in
$\varepsilon_r$.  For~(i), $\varepsilon_r<0$ forces $\mu_p>1-r>0$ and hence
$\rho_c>0$, so $0<\Lambda<1$.
Result~\ref{res:negative-pt} translates the sign of $\rho_c$ into conditions
on $\mu_p$ and $|\varepsilon_S|\gtrless -\varepsilon_D$: with
$\varepsilon_r>0$ this is~(ii), and $\rho_c<0$ is~(iii).
\end{proof}
\begin{proof}[Proof of Corollary~\ref{cor:welfare-bias-buyer}]
The true change is $-Q\mu_p\Lambda$ and the no-return prediction is $-Q$, so
the no-return loss figure $Q$ understates the returns-adjusted figure
$Q\mu_p\Lambda$ if and only if $\mu_p\Lambda>1$ and overstates it if and only
if $\mu_p\Lambda<1$.  Write
$\mu_p=1-r(p)-r'(p)(p+h)$ and $\Lambda=1+r'(p)(p+h)\rho_c$.  Using
$\mu_p\rho_c=\varepsilon_S/(\varepsilon_S-\varepsilon_D)$,
\begin{align*}
  \mu_p\Lambda-1
  &=\mu_p\bigl(1+r'(p)(p+h)\rho_c\bigr)-1
  \\
  &=\mu_p-1+r'(p)(p+h)\,(\mu_p\rho_c)
  \\
  &=-r(p)-r'(p)(p+h)
    +r'(p)(p+h)\,\frac{\varepsilon_S}{\varepsilon_S-\varepsilon_D}
  \\
  &=-r(p)+r'(p)(p+h)\,\frac{\varepsilon_D}{\varepsilon_S-\varepsilon_D}
  \\
  &=-r(p)\Biggl(1-\varepsilon_r(p)\,\frac{p+h}{p}\,
       \frac{\varepsilon_D}{\varepsilon_S-\varepsilon_D}\Biggr).
\end{align*}
Hence $\mu_p\Lambda<1$ if and only if
$\varepsilon_r(p)\,\varepsilon_D/(\varepsilon_S-\varepsilon_D)<p/(p+h)$.
Multiplying by $-1$ gives
\[
  \varepsilon_r(p)\,
  \frac{-\varepsilon_D}{\varepsilon_S-\varepsilon_D}
  >
  -\frac{p}{p+h},
\]
which is \eqref{eq:V-buyer-cut}.  The reverse strict inequality is
$\mu_p\Lambda>1$.
\end{proof}

\section{Tax partial of the return rate}
\label{app:rtau}

The average return rate is $\tilde{r}(p,\tau)=N/D$ with
\[
  N=\int_p^\infty \phi(v,p-\tau)\,dF(v;\tau),
  \qquad
  D=1-F(p;\tau).
\]
Differentiating at fixed $p$,
\[
  \tilde{r}_{\tau}
  =\frac{N_{\tau}D-ND_{\tau}}{D^{2}}
  =\frac{N_{\tau}}{D}-\tilde{r}\,\frac{D_{\tau}}{D}.
\]
At $\tau=0$, $D=D(p)$ and $\tilde{r}=r(p)$.  Among remaining buyers, recovery
falls one-for-one, which contributes
$-\int_p^\infty \partial\phi(v,p)/\partial p\,dF(v)$ to $N_{\tau}$.  Buyer mass
changes at rate $D_{\tau}\le 0$.  The types who leave are those at $v=p$.  Over
an interval $d\tau$ they have measure $-D_{\tau}\,d\tau$ and cell return
probability $\phi(p,p)$, so $N$ changes by
\[
  -\phi(p,p)\,(-D_{\tau})\,d\tau
  =\phi(p,p)\,D_{\tau}\,d\tau.
\]
Collecting the two pieces,
\[
  N_{\tau}
  =-\int_p^\infty \frac{\partial\phi(v,p)}{\partial p}\,dF(v)
  +\phi(p,p)\,D_{\tau}.
\]
Substituting yields
\begin{equation}
  \tilde{r}_{\tau}
  =-\frac{1}{D(p)}\int_p^\infty \frac{\partial\phi(v,p)}{\partial p}\,dF(v)
  +\frac{-D_{\tau}}{D(p)}\bigl(r(p)-\phi(p,p)\bigr).
  \label{eq:rtau-split}
\end{equation}


\bibliographystyle{plainnat}
\bibliography{references}

\end{document}
